% GENERATED FILE. Edit canonical lessons, then run npm run generate:books.
% canonical-source-hash: fnv1a64:3e710befbad3b61d
% canonical-lessons: PA-C181-rel, PA-C181-ticket, PA-C181-yayya, PA-C181-yatra, PA-C181-vapas, PA-C181-paidal, PA-C181-travel-review, PA-C181-payoff

\chapter{A Train Ticket and the Way Home}
\label{ch:pa-a2-transport-prerequisite}
\hlchaptermodality{181}

\begin{chapteropening}
\textbf{By the end of this chapter:} \emph{I can understand and give one short practical note about a train ticket, a journey, and going home on foot.}

Recognise the newly needed \pa{ਯ}, then retrieve a three-line travel note through listening, speaking, reading, and bounded Gurmukhi writing after an atom-free review.
\end{chapteropening}

\section[\texorpdfstring{\pa{ਰੇਲ} (rel)}{ਰੇਲ (rel)}]{\pa{ਰੇਲ} — rail or train}

\label{lesson:PA-C181-rel}



\begin{quote}
Recall \textbf{\pa{ਸਟੇਸ਼ਨ}}, then the older train word \textbf{\pa{ਗੱਡੀ}}.
\end{quote}

\begin{grammarlens}[title={Grammar Lens: one transport word}]
\begin{quote}
\textbf{\pa{ਰੇਲ}} — \emph{rel} — rail; train
\end{quote}

Keep this as one practical transport word. The older \textbf{\pa{ਗੱਡੀ}} remains useful.
\end{grammarlens}

\begin{cousinweb}
\textbf{\pa{ਰੇਲ}} (\emph{rel}) is a loan from English \emph{rail}.
\end{cousinweb}

\subsection*{Your turn}
\begin{itemize}
\raggedright
  \item \emph{Hear:} \emph{rel}; YOU SAY: \textbf{\pa{ਰੇਲ}}
  \item \emph{Read it:} \textbf{\pa{ਰੇਲ}}
  \item \emph{Say it:} \textbf{\pa{ਸਟੇਸ਼ਨ}}, \textbf{\pa{ਗੱਡੀ}}, \textbf{\pa{ਰੇਲ}}
\end{itemize}

\begin{tcolorbox}[breakable,colback=teal!4,colframe=teal!35!black,title={Before you move on}]
Say \emph{rail; train}. (\textbf{\pa{ਰੇਲ}}.)

Source: Punjabi University, \href{https://dic.learnpunjabi.org/default.aspx?look=\%E0\%A8\%B0\%E0\%A9\%87\%E0\%A8\%B2}{RCPLT dictionary}, lists \textbf{\pa{ਰੇਲ}} as rail or train.
\end{tcolorbox}
\section[\texorpdfstring{\pa{ਟਿਕਟ} (tikat)}{ਟਿਕਟ (tikat)}]{\pa{ਟਿਕਟ} — ticket}

\label{lesson:PA-C181-ticket}



\begin{quote}
Say the new rail word once. (\textbf{\pa{ਰੇਲ}}.)
\end{quote}

\begin{grammarlens}[title={Grammar Lens: one travel item}]
\begin{quote}
\textbf{\pa{ਟਿਕਟ}} — \emph{ṭikaṭ} — ticket
\end{quote}

\begin{quote}
\textbf{\pa{ਰੇਲ} \pa{ਦੀ} \pa{ਟਿਕਟ}} — \emph{rel dī ṭikaṭ} — a train ticket
\end{quote}

Use the whole short phrase first; it adds no new possession table.
\end{grammarlens}

\begin{cousinweb}
\textbf{\pa{ਟਿਕਟ}} (\emph{ṭikaṭ}) is a loan from English \emph{ticket}.
\end{cousinweb}

\subsection*{Your turn}
\begin{itemize}
\raggedright
  \item \emph{Hear:} \emph{ṭikaṭ}; YOU SAY: \textbf{\pa{ਟਿਕਟ}}
  \item \emph{Read it:} \textbf{\pa{ਰੇਲ} \pa{ਦੀ} \pa{ਟਿਕਟ}}
\end{itemize}

\begin{tcolorbox}[breakable,colback=teal!4,colframe=teal!35!black,title={Before you move on}]
Say \emph{train ticket}. (\textbf{\pa{ਰੇਲ} \pa{ਦੀ} \pa{ਟਿਕਟ}}.)

Source: Punjabi University, \href{https://dic.learnpunjabi.org/default.aspx?look=\%E0\%A8\%9F\%E0\%A8\%BF\%E0\%A8\%95\%E0\%A8\%9F}{RCPLT dictionary}, lists \textbf{\pa{ਟਿਕਟ}} as a ticket.
\end{tcolorbox}
\section[\texorpdfstring{\pa{ਯ} (ya)}{ਯ (ya)}]{\pa{ਯ} — one new letter before a journey word}

\label{lesson:PA-C181-yayya}



\begin{quote}
Say \textbf{\pa{ਰੇਲ} \pa{ਦੀ} \pa{ਟਿਕਟ}} once.
\end{quote}

\subsection*{Script you'll notice: \pa{ਯ}}
\begin{quote}
\pa{ਯ}
\end{quote}

\textbf{yayyā} — \emph{ya}. Keep the letter in view. Today you only need to recognise it and connect it with \emph{ya}.

No stroke order is taught here: the track does not yet have a verified source for one.

\subsection*{Writing: copy what you can see}
Keep \textbf{\pa{ਯ}} in front of you. Trace its printed shape once with a finger, then copy it once with a pencil while the model stays visible. Say \emph{ya} when you finish.

\begin{quote}
This is observe-and-copy with the model visible, not writing from memory.
\end{quote}

\subsection*{Your turn}
\begin{itemize}
\raggedright
  \item \emph{Look:} at \textbf{\pa{ਯ}}; YOU SAY: \emph{ya}
  \item Look away. Then look back and point to \textbf{\pa{ਯ}}.
  \item \emph{Say it:} \textbf{\pa{ਰੇਲ} \pa{ਦੀ} \pa{ਟਿਕਟ}}, then \emph{ya}
\end{itemize}

\begin{tcolorbox}[breakable,colback=teal!4,colframe=teal!35!black,title={Before you move on}]
Point to \textbf{\pa{ਯ}} and say its sound. (\emph{ya}.)

Source: \href{https://www.unicode.org/charts/PDF/U0A00.pdf}{Unicode Gurmukhi chart}.
\end{tcolorbox}
\section[\texorpdfstring{\pa{ਯਾਤਰਾ} (yatra)}{ਯਾਤਰਾ (yatra)}]{\pa{ਯਾਤਰਾ} — journey or travel}

\label{lesson:PA-C181-yatra}



\begin{quote}
Retrieve \textbf{\pa{ਰੇਲ} \pa{ਦੀ} \pa{ਟਿਕਟ}}, then point to \textbf{\pa{ਯ}} and say \emph{ya}.
\end{quote}

\begin{grammarlens}[title={Grammar Lens: one journey word}]
\begin{quote}
\textbf{\pa{ਯਾਤਰਾ}} — \emph{yātrā} — journey; travel
\end{quote}

\begin{quote}
\textbf{\pa{ਰੇਲ} \pa{ਨਾਲ} \pa{ਯਾਤਰਾ}} — \emph{rel nāl yātrā} — travel by train
\end{quote}

Treat the second line as one useful travel chunk.
\end{grammarlens}

\begin{cousinweb}
\textbf{\pa{ਯਾਤਰਾ}} (\emph{yātrā}) is Sanskrit \emph{yātrā}, “journey, procession,” taken over as a learned loan.
\end{cousinweb}

\subsection*{Your turn}
\begin{itemize}
\raggedright
  \item \emph{Hear:} \emph{yātrā}; YOU SAY: \textbf{\pa{ਯਾਤਰਾ}}
  \item \emph{Read it:} \textbf{\pa{ਰੇਲ} \pa{ਨਾਲ} \pa{ਯਾਤਰਾ}}
  \item \emph{Say it:} \textbf{\pa{ਰੇਲ} \pa{ਦੀ} \pa{ਟਿਕਟ}}, then \textbf{\pa{ਰੇਲ} \pa{ਨਾਲ} \pa{ਯਾਤਰਾ}}
\end{itemize}

\begin{tcolorbox}[breakable,colback=teal!4,colframe=teal!35!black,title={Before you move on}]
Say \emph{travel by train}. (\textbf{\pa{ਰੇਲ} \pa{ਨਾਲ} \pa{ਯਾਤਰਾ}}.)

Source: Punjabi University, \href{https://dic.learnpunjabi.org/default.aspx?look=\%E0\%A8\%AF\%E0\%A8\%BE\%E0\%A8\%A4\%E0\%A8\%B0\%E0\%A8\%BE}{RCPLT dictionary}, lists \textbf{\pa{ਯਾਤਰਾ}} as journey or travel.
\end{tcolorbox}
\section[\texorpdfstring{\pa{ਵਾਪਸ} (vapas)}{ਵਾਪਸ (vapas)}]{\pa{ਵਾਪਸ} — back or returning}

\label{lesson:PA-C181-vapas}



\begin{quote}
Say \emph{journey}, then \emph{home}. (\textbf{\pa{ਯਾਤਰਾ}}, \textbf{\pa{ਘਰ}}.)
\end{quote}

\begin{grammarlens}[title={Grammar Lens: one direction word}]
\begin{quote}
\textbf{\pa{ਵਾਪਸ}} — \emph{vāpas} — back; returning
\end{quote}

\begin{quote}
\textbf{\pa{ਵਾਪਸ} \pa{ਘਰ}} — \emph{vāpas ghar} — back home
\end{quote}

Keep this as one direction word and one short chunk.
\end{grammarlens}

\subsection*{Your turn}
\begin{itemize}
\raggedright
  \item \emph{Hear:} \emph{vāpas}; YOU SAY: \textbf{\pa{ਵਾਪਸ}}
  \item \emph{Read it:} \textbf{\pa{ਵਾਪਸ} \pa{ਘਰ}}
  \item \emph{Say it:} \textbf{\pa{ਯਾਤਰਾ}}, then \textbf{\pa{ਵਾਪਸ} \pa{ਘਰ}}
\end{itemize}

\begin{tcolorbox}[breakable,colback=teal!4,colframe=teal!35!black,title={Before you move on}]
Say \emph{back home}. (\textbf{\pa{ਵਾਪਸ} \pa{ਘਰ}}.)

Source: Punjabi University, \href{https://dic.learnpunjabi.org/default.aspx?look=\%E0\%A8\%B5\%E0\%A8\%BE\%E0\%A8\%AA\%E0\%A8\%B8}{RCPLT dictionary}, lists \textbf{\pa{ਵਾਪਸ}} as back or returning.
\end{tcolorbox}
\section[\texorpdfstring{\pa{ਪੈਦਲ} (paidal)}{ਪੈਦਲ (paidal)}]{\pa{ਪੈਦਲ} — on foot}

\label{lesson:PA-C181-paidal}



\begin{quote}
Retrieve \textbf{\pa{ਵਾਪਸ} \pa{ਘਰ}}, then the old verb \textbf{\pa{ਜਾਣਾ}}.
\end{quote}

\begin{grammarlens}[title={Grammar Lens: one way of travelling}]
\begin{quote}
\textbf{\pa{ਪੈਦਲ}} — \emph{paidal} — on foot
\end{quote}

\begin{quote}
\textbf{\pa{ਪੈਦਲ} \pa{ਵਾਪਸ} \pa{ਘਰ} \pa{ਜਾਣਾ}} — \emph{paidal vāpas ghar jāṇā} — to go back home on foot
\end{quote}

Use this as one bounded travel chunk. No new verb ending appears here.
\end{grammarlens}

\subsection*{Your turn}
\begin{itemize}
\raggedright
  \item \emph{Hear:} \emph{paidal}; YOU SAY: \textbf{\pa{ਪੈਦਲ}}
  \item \emph{Read it:} \textbf{\pa{ਪੈਦਲ} \pa{ਵਾਪਸ} \pa{ਘਰ} \pa{ਜਾਣਾ}}
  \item \emph{Say it:} \textbf{\pa{ਵਾਪਸ} \pa{ਘਰ}}, then rebuild the full chunk
\end{itemize}

\begin{tcolorbox}[breakable,colback=teal!4,colframe=teal!35!black,title={Before you move on}]
Say \emph{to go back home on foot}. (\textbf{\pa{ਪੈਦਲ} \pa{ਵਾਪਸ} \pa{ਘਰ} \pa{ਜਾਣਾ}}.)

Source: Punjabi University, \href{https://dic.learnpunjabi.org/default.aspx?look=\%E0\%A8\%AA\%E0\%A9\%88\%E0\%A8\%A6\%E0\%A8\%B2}{RCPLT dictionary}, lists \textbf{\pa{ਪੈਦਲ}} as on foot.
\end{tcolorbox}
\section[\texorpdfstring{\pa{ਰੇਲ} \pa{ਦੀ} \pa{ਟਿਕਟ} · \pa{ਰੇਲ} … (rel di tikat …)}{ਰੇਲ ਦੀ ਟਿਕਟ · ਰੇਲ … (rel di tikat …)}]{Three small travel chunks — guided review}

\label{lesson:PA-C181-travel-review}



\begin{quote}
Recall \textbf{\pa{ਰੇਲ}}, \textbf{\pa{ਟਿਕਟ}}, \textbf{\pa{ਯਾਤਰਾ}}, \textbf{\pa{ਵਾਪਸ}}, and \textbf{\pa{ਪੈਦਲ}}.
\end{quote}

\subsection*{Your turn: fit the known pieces}
\begin{quote}
\textbf{\pa{ਰੇਲ} \pa{ਦੀ} \pa{ਟਿਕਟ}} — a train ticket
\end{quote}

\begin{quote}
\textbf{\pa{ਰੇਲ} \pa{ਨਾਲ} \pa{ਯਾਤਰਾ}} — travel by train
\end{quote}

\begin{quote}
\textbf{\pa{ਪੈਦਲ} \pa{ਵਾਪਸ} \pa{ਘਰ}} — back home on foot
\end{quote}

\begin{itemize}
\raggedright
  \item \emph{Hear:} each chunk; YOU SAY: it once
  \item \emph{Read it:} all three chunks
  \item \emph{Say it:} the ticket chunk, the journey chunk, then the homeward chunk
\end{itemize}

\begin{tcolorbox}[breakable,colback=teal!4,colframe=teal!35!black,title={Before you move on}]
Rebuild all three travel chunks without looking.

Sources: Punjabi University RCPLT entries for \href{https://dic.learnpunjabi.org/default.aspx?look=\%E0\%A8\%B0\%E0\%A9\%87\%E0\%A8\%B2}{\pa{ਰੇਲ}}, \href{https://dic.learnpunjabi.org/default.aspx?look=\%E0\%A8\%9F\%E0\%A8\%BF\%E0\%A8\%95\%E0\%A8\%9F}{\pa{ਟਿਕਟ}}, \href{https://dic.learnpunjabi.org/default.aspx?look=\%E0\%A8\%AF\%E0\%A8\%BE\%E0\%A8\%A4\%E0\%A8\%B0\%E0\%A8\%BE}{\pa{ਯਾਤਰਾ}}, \href{https://dic.learnpunjabi.org/default.aspx?look=\%E0\%A8\%B5\%E0\%A8\%BE\%E0\%A8\%AA\%E0\%A8\%B8}{\pa{ਵਾਪਸ}}, and \href{https://dic.learnpunjabi.org/default.aspx?look=\%E0\%A8\%AA\%E0\%A9\%88\%E0\%A8\%A6\%E0\%A8\%B2}{\pa{ਪੈਦਲ}}.
\end{tcolorbox}
\section[\texorpdfstring{\pa{ਰੇਲ} \pa{ਦੀ} \pa{ਟਿਕਟ} \pa{ਇੱਥੇ} \pa{ਹੈ।} … (rel di tikat …)}{ਰੇਲ ਦੀ ਟਿਕਟ ਇੱਥੇ ਹੈ। … (rel di tikat …)}]{One bounded travel note}

\label{lesson:PA-C181-payoff}



\begin{quote}
Say the three travel chunks from the review.
\end{quote}

\subsection*{Your turn: four-skill payoff}
\begin{quote}
\textbf{\pa{ਰੇਲ} \pa{ਦੀ} \pa{ਟਿਕਟ} \pa{ਇੱਥੇ} \pa{ਹੈ।} \pa{ਮੈਨੂੰ} \pa{ਰੇਲ} \pa{ਨਾਲ} \pa{ਯਾਤਰਾ} \pa{ਪਸੰਦ} \pa{ਹੈ।} \pa{ਮੈਨੂੰ} \pa{ਪੈਦਲ} \pa{ਵਾਪਸ} \pa{ਘਰ} \pa{ਜਾਣਾ} \pa{ਪਸੰਦ} \pa{ਹੈ।}}
\end{quote}

\begin{itemize}
\raggedright
  \item \emph{Hear:} the three-line travel note; YOU SAY: it one line at a time
  \item \emph{Read it:} the whole travel note
  \item \emph{Write it:} the three-line travel note after covering the model
\end{itemize}

\begin{tcolorbox}[breakable,colback=teal!4,colframe=teal!35!black,title={Before you move on}]
Rebuild the ticket line, journey line, and homeward line without looking.

This is the fourth prerequisite tranche exposed by Punjabi A2 mock drafting. It removes 22 measured unsupported occurrences but does not unblock Punjabi A2 mock 1 by itself. It is not an A2 mock, not A2 readiness or attainment, and not external or human validation.

Sources: Punjabi University RCPLT entries for \href{https://dic.learnpunjabi.org/default.aspx?look=\%E0\%A8\%B0\%E0\%A9\%87\%E0\%A8\%B2}{\pa{ਰੇਲ}}, \href{https://dic.learnpunjabi.org/default.aspx?look=\%E0\%A8\%9F\%E0\%A8\%BF\%E0\%A8\%95\%E0\%A8\%9F}{\pa{ਟਿਕਟ}}, \href{https://dic.learnpunjabi.org/default.aspx?look=\%E0\%A8\%AF\%E0\%A8\%BE\%E0\%A8\%A4\%E0\%A8\%B0\%E0\%A8\%BE}{\pa{ਯਾਤਰਾ}}, \href{https://dic.learnpunjabi.org/default.aspx?look=\%E0\%A8\%B5\%E0\%A8\%BE\%E0\%A8\%AA\%E0\%A8\%B8}{\pa{ਵਾਪਸ}}, and \href{https://dic.learnpunjabi.org/default.aspx?look=\%E0\%A8\%AA\%E0\%A9\%88\%E0\%A8\%A6\%E0\%A8\%B2}{\pa{ਪੈਦਲ}}.
\end{tcolorbox}
